\section{Low modes: geometry and cell states}
\label{sec:low-modes-geometry-and-cell-states}

It remains to treat the bounded, low direct-bias and low cluster modes,
one mode at a time. We first assign all odd roles within small cells,
reserving some of these roles for replacement later. A cell keeps a fixed
pool of host bins. One \emph{fresh cell state} consists of its primitive
slice data, its chosen group bins, its injective odd labels, and the
internal even priors computed from those data. We will construct a sampler
whose bin and label stages preserve their specified conditional singleton
marginals and satisfy upper comparisons for joint observations.
Section~\ref{sec:initial-list-gates-and-resampling-estimates}
uses independent repetitions of this sampler to prepare the early lists.

Recall \(C_n=N2^{-n}\). In direct modes take singleton bins \(d_i=1\)
and singleton groups on odd roles. Let \(\mathcal X_i\) be the cleaned
envelope of patch \(i\), which is the single cleaned support in direct
modes. The scale bounds give, uniformly in the present cases,
\[
\begin{aligned}
 h_i&\le(\log n)^{.1},& \ell_i&=o(\log n),\\
 M_i&=N n^{-o(1)},& d_i&=n^{o(1)},\\
 n|D_i-.5|&=O(\log n)\quad\text{on }\mathcal X_i,&s_i&=O(\log n).
\end{aligned}
\]
The constants may depend on the fixed large thresholds. In low cluster
mode, neither \(d_i\) nor \(h_i\) need tend to infinity; the choice of
\(Q_0\) ensures the required bin and internal-parameter inequalities.
Call \(x,z\) conflicting for patch \(i\) when
\(|K_{\pi_i}(x,z)|>\xi\), including the diagonal. By
Lemma~\ref{lem:homogeneous-peeling} and the cleaning, the number of labels
conflicting with a fixed \(x\) in any one used corner support is at most
\(\Delta_i\le K e^{s_i}\). This assertion concerns each corner support
separately. Fixed factors may be incorporated into \(K\) below.

\begin{lemma}[Late classes and cell inputs]\label{lem:late-classes-cells}
The low-mode cube admits the late-class and cell allocations described
below, with \(A_0\log n\le r<2A_0\log n\) late classes. For
\(0\le j\le r\), every even star has between
\(\lfloor j/2\rfloor\) and \(j\) neighbors in the last \(j\) classes,
at most one in each class, and at most one late internal neighbor.
Distinct slices of a cell have the separation stated below. Pool inputs
are uniform permutation slots; conditional on the pools, fresh cell states
can be drawn using independent local tapes.
\end{lemma}

\Needspace{5\baselineskip}
\begin{proof}\leavevmode\par
\paragraph{Syndrome classes.}
Let \(\mathcal H\) be a binary group of power-of-two size
\(n\le H_*<2n\). Give the \(n\) coordinates distinct IDs
\(a_1,\ldots,a_n\in\mathcal H\), including every element of a linear
hyperplane \(\mathcal H_0\). Choose a subspace \(\mathcal L\) with
\(A_0\log n\le r:=|\mathcal L|<2A_0\log n\), not contained in
\(\mathcal H_0\). The IDs of all axes in the nested internal allocation
can be put in distinct cosets modulo \(\mathcal L\). There are enough
such cosets even using IDs in \(\mathcal H_0\), since the total number of
internal axes is at most \((\log n)^{.1}\), whereas
\(|\mathcal H/\mathcal L|\asymp n/\log n\).

For a cube word \(z\), set \(s(z)=\sum_k a_kz_k\). An odd role is late
when \(s(z)\in\mathcal L\); the \(r\) possible syndromes give its late
classes. Each class has size \(2^{n-1}/H_*\): the joint map recording
syndrome and parity has full rank, since the IDs include zero and span
\(\mathcal H\). Order the classes so that every suffix divides as evenly
as possible between syndromes in and outside \(\mathcal H_0\).

At an even \(v\), flipping coordinate \(k\) reaches class \(t\) exactly
when \(a_k=s(v)+t\). If \(s(v)\in\mathcal H_0\), every
\(t\in\mathcal L\cap\mathcal H_0\) has its required ID in
\(\mathcal H_0\); if \(s(v)\notin\mathcal H_0\), the same is true of
every \(t\in\mathcal L\setminus\mathcal H_0\). Thus any last \(j\)
classes supply at least \(\lfloor j/2\rfloor\) neighbors. Distinct IDs
allow at most one per class, giving the upper bound \(j\). Write
\(p(v)\) for the total number of late neighbors, so \(r/2\le p(v)\le r\).
All eligible flip IDs lie in the single coset \(s(v)+\mathcal L\), and
hence at most one is internal.

Split the reserved \(Y\)-labels into disjoint late pools of size
\(\asymp N/r\). The initial assignment nevertheless labels every odd
role, using only the \(Y_i\)'s outside this reserve. The initial labels
of late roles are dummies, to be replaced later. The remaining odd roles
are called early. The initial lists will impose only the early external
hits; an internal prior may still retain the dummy hit at the one late
internal neighbor. Deferring the other late hits saves a factor
\(2^{-p(v)}\), up to a fixed constant, in the initial atom estimate
of Lemma~\ref{lem:independent-pinned-list}. This supplies polynomially
small initial failure bounds even when the patch gain \(s_i\) is bounded.

\paragraph{Separated cells and persistent pools.}
Partition each patch into cells of at most \(n^{A_c}\) positions, each a
union of entire slices, so that distinct slices of one cell have cube
distance greater than \((\log n)^3\). To do this, greedily color the
slices with \(\exp(O(\log^4 n))\) colors so that near slices have
different colors, then batch each color. This gives
\(O(2^{n-\ell_i}/n^{A_c})\) cells in patch \(i\).

Give each cell
\[
 L_i^{\rm slot}=\lceil K_{\rm cell}n^{A_c}/d_i\rceil
              =n^{A_c+o(1)}
\]
slots, where \(K_{\rm cell}\) is a sufficiently large fixed constant
compared with \(1/\theta_*\). The slot count is distinct from all
neighbor counts. There are \(B_i=M_i/d_i\) physical bins in the patch.
The patch counts and \(C_n\to\infty\) ensure that all cells' slots fit
disjointly among these bins. A uniform permutation of the bins, independent
for each patch, supplies the slot images and thus the cell pools.

On any \(\exp(\mathrm{polylog}(n))\) specified slots, with fixed powers
of \(\log n\), this law has upper comparison \(1+o(1)\) with iid
uniform bins. The same comparison holds after fixing any one global slot
to a prescribed bin, retaining that pin in both laws. Indeed the
without-replacement factors differ from one by a total
\(\exp(\mathrm{polylog}(n))/B_i=o(1)\), since \(B_i=2^{n+o(n)}\).

The pool persists when a cell is resampled. Its tape is a sequence of
independent randomizers for the conditional fresh-state sampler constructed
below. A drawing reads only this cell's pool and the selected tape entry;
it uses local fallback rules on invalid inputs. Later schedules may depend
on other data when deciding which entry to read.
\end{proof}

\paragraph{A fixed table of permitted bins.}
We need a label in an externally prescribed slot to be compatible with a
fresh internal prior. For an external early incidence \((b,v)\), meaning
that the edge is noninternal, remove from the entire group of \(b\) every
bin containing a label \(y\) for which
\begin{equation}
 \Pr_v^{\rm base}\{\sigma_v\ne0,\ |d_G(\sigma_v,y)-.5|>2b_*\}
       >e^{-c n}.                                                \label{eq:source-21}
\end{equation}
Here \(\Pr_v^{\rm base}\) is the unrestricted internal \(W\) experiment
and reference label sampling at \(v\), or the deterministic prior in
direct modes. It is an independent comparison experiment, not the eventual
cell state at \(v\). The constant \(c>0\) is sufficiently small.

Each nonzero \(\sigma_v\) has small width by \eqref{eq:source-14} or
the direct bounds. Deep discrepancy therefore bounds the probability in
\eqref{eq:source-21}, averaged over uniform \(Y_{i(b)}\), by
\(e^{-3cn}\). Markov, followed by multiplication by the bin size and
the number of incidences of a group, shows that an \(e^{-\Omega(n)}\)
fraction of all bins is removed. The remaining bins form a deterministic
permission table, fixed before the cells are sampled. Write
\(\mathcal B_g^{\rm perm}\) for the bins permitted at group \(g\).

In cluster mode begin with the internally pretrimmed bin law
\(q_g^{\mathrm{in},W}\) used to define the low output \(\pi_i\) in
Proposition~\ref{prop:balanced-patch-profiles}. Its further restriction to
permitted bins is
\[
 \bar q_g(D)=
 \frac{q_g^{\mathrm{in},W}(D)\mathbf1_{\{D\in\mathcal B_g^{\rm perm}\}}}
      {q_g^{\mathrm{in},W}(\mathcal B_g^{\rm perm})}.
\]
This additional loss is \(e^{-\Omega(n)}\), by \eqref{eq:source-13}.
The in-bin law \(U_g^W(D)\) is unchanged. The same notation applies in
direct modes, starting from uniform singleton bins and using point in-bin
laws. Although the permission table is fixed, \(\bar q_g\) still depends
on \(W\) in cluster mode. Throughout these changes of sampling law,
\(\sigma_v\) and \(E_v^{\rm in}\) remain the raw posterior and failure
rules of Section~\ref{sec:internal-even-posteriors}.

\subsection{Pool and history conditions}
\label{sec:pool-and-history-conditions}

Fix one cell. In cluster mode first draw its slice histories independently,
each conditioned on \eqref{eq:source-15} throughout that slice. Given
these histories and the cell pool \(\mathcal P\), define, whenever the
denominator is positive,
\[
 \widetilde q_g(D)=
 \frac{\bar q_g(D)\mathbf1_{\{D\in\mathcal P\}}}
      {\sum_{D'\in\mathcal P}\bar q_g(D')}.
\]
A zero denominator uses the local fallback. The same restriction defines
the direct-mode laws without \(W\).
A pool is individually typical when the following hold:
\begin{itemize}
\item Its slots have distinct images, and every denominator above is
\((L_i^{\rm slot}/B_i)(1+O(n^{-4}))\), with relative error at most
\(n^{-4}\), uniformly over every group and every admissible history of
its slice.
\item In cluster mode, for every such history and internal star \(v\),
independent group choices from \(\widetilde q_g\), followed by independent
reference role labels given those bins, have
\(\Pr(E_v^{\rm in})\le\epsilon_i^{1/4}\). This also holds if one
participating group's bin is fixed to any target of positive restricted
mass.
\end{itemize}
The latter is a \emph{group-bin pin}. A \emph{slot pin} instead fixes one
pool slot to a physical bin. Both are allowed simultaneously in the
estimates that follow. Later a \emph{role-label pin} will fix an output
of the within-bin injection.

\begin{lemma}[Typical pools and the history-load gate]\label{lem:typical-cell-pools}
The preceding pool requirements fail at a cell with probability at most
\(e^{-n^{c_0}}\), for a fixed \(c_0>0\), also conditional on any one
prescribed slot pin. On a typical pool the additional history-load gate
below has negligible relative conditioning cost. Direct modes already
have the required load slack.
\end{lemma}

\Needspace{5\baselineskip}
\begin{proof}\leavevmode\par
\paragraph{Empirical normalizers.}
Use the iid-slot comparison, keeping a prescribed slot pin if present.
For fixed admissible slice data and slots \(D_1,\ldots,D_{L_i^{\rm slot}}\),
write
\[
 \widehat q_g=
 \frac{B_i}{L_i^{\rm slot}}\sum_l\bar q_g(D_l)\delta_{D_l}.
\]
Its total mass is the pool restriction denominator in units of
\(L_i^{\rm slot}/B_i\). Changing one slot changes this mass by at most
\(O(\exp(O(k_iT_i))/L_i^{\rm slot})\); in direct modes the exponential
factor is constant. Its expectation is one, up to the negligible change
caused by a slot pin. Bounded differences gives the required relative
\(n^{-4}\) accuracy.

\paragraph{Empirical failure probabilities.}
Let \(g_1,\ldots,g_m\) be the participating groups at an internal star,
and set
\[
 f_v^W(D_1',\ldots,D_m')
   =\Pr_{\rm independent\ roles}(E_v^{\rm in}\mid W,D_1',\ldots,D_m').
\]
Before pool restriction, its integral against the \(\bar q_g\)'s is at
most \(2\sqrt{\epsilon_i}\), including after a positive group-bin pin.
This follows from the internal pretrim and \eqref{eq:source-15}; the
normalization losses over the other groups fit because
\(h_i^3\sqrt{\epsilon_i}\ll1\), and the permission trim is exponentially
small.

Consider instead the unnormalized integral
\[
 \int f_v^W\,d\widehat q_{g_1}\cdots d\widehat q_{g_m},
\]
replacing one factor by a point mass if a group-bin pin is imposed. Its
change under one slot replacement is \(n^{-A_c+o(1)}\), since
\(\exp(O(h_i k_iT_i))=n^{o(1)}\). To compute its expectation, expand the
empirical measures by their slot indices. Distinct unpinned indices give
back the original independent bin laws. Coincident indices, or occurrences
of a pinned slot, contribute at most \(n^{-A_c+o(1)}\). Thus its
expectation is at most \(2\sqrt{\epsilon_i}+n^{-A_c+o(1)}\).

Bounded differences puts this integral below
\(\epsilon_i^{1/4}/2\) except with exponentially small probability;
normalizing the empirical masses then gives the claimed bound.
Here \(\epsilon_i=n^{-o(1)}\), and
\(2\sqrt{\epsilon_i}\ll\epsilon_i^{1/4}\), uniformly for the fixed
large threshold. Both concentration exponents can exceed \(n^{100}\).
By Lemma~\ref{lem:finite-low-mode-data}, the logarithm of the number of
slice histories is at most \(n^{1+o(1)}\); adding group-bin targets costs
only \(O(n)\), and there are polynomially many slices in a cell. The
union therefore preserves an \(e^{-n^{c_0}}\) bound. Repeated slot
images have probability \(e^{-\Omega(n)}\), including after one slot pin.

\paragraph{The load gate on slice histories.}
For a fixed typical pool in cluster mode, let
\[
 \Lambda_y(W)=\sum_{b\ {\rm in\ the\ cell}}
     \widetilde q_{g(b)}(D(y))U_{g(b)}^W(D(y))(y)
\]
be the conditional expected load on label \(y\). Condition the slice
histories further on \(\Lambda_y(W)\le\theta_*\) for every label.
To bound the cost, normalizer control and the permission-trim bound compare
each role's law with \((1+O(n^{-3}))(B_i/L_i^{\rm slot})\) times its
internally pretrimmed law. Averaging the latter gives \(\pi_i\), so
\eqref{eq:source-16} bounds the role's mean at a pool label by
\(12/(L_i^{\rm slot}d_i)\). Hence
\(\mathbb E\Lambda_y=O(1/K_{\rm cell})\).

Different slices are independent under the slice-conditioned law. Each
slice contributes at most \(n^{-A_c+o(1)}\) to a column, by
\eqref{eq:source-13} and the restriction normalizers. Concentration and
a union over labels show that the load gate fails with probability
\(\delta(\mathcal P)\le e^{-n^{c_0}}\), decreasing \(c_0\) if necessary.
For every nonnegative test \(F\), consequently,
\[
 \mathbb E[F\mid\mathcal P,\ \text{load gate}]
 \le\frac{1}{1-\delta(\mathcal P)}\,
       \mathbb E[F\mid\mathcal P].
\]
The expectations here already use slice-conditioned histories. This
conditioning is performed separately at each fixed pool; it does not
reweight the pool law. In direct modes typicality itself gives the same
load slack, with no history conditioning.
\end{proof}

\subsection{Cell calibrations}
\label{sec:cell-calibrations}

We now complete the sampler at a typical pool and a passing history.
The two calibration steps keep exact single-coordinate marginals while
imposing the internal tests and distinct labels.

\begin{proposition}[Calibrated group bins]\label{prop:calibrated-group-bins}
Given those successful data in a cluster cell, there is a bin sampler
across groups with exact individual marginals \(\widetilde q_g\) and the
following properties:
\begin{itemize}
\item At every internal star, different participating groups choose
distinct bins, and the independent-role conditional probability of
\(E_v^{\rm in}\) given bins is at most \(\epsilon_i^{1/8}\).
\item Given-bin label column sums are at most \(.2\).
\item For targets on \(l\) distinct groups, the joint probability is at
most \(\exp(d_i^{-.05}l)\) times their product \(\widetilde q_g\)-mass.
\end{itemize}
\end{proposition}

\begin{proof}\leavevmode\par
Fix the data throughout. Let \(\mathcal C\) be the set of laws on bin
assignments that satisfy the first two requirements and every joint upper
inequality in the third, always relative to the original
\(\widetilde q_g\)'s. This is a compact convex set. We construct members
of \(\mathcal C\) for arbitrary small price-directed perturbations, then
separate its marginal image to obtain exact calibration.

\paragraph{Avoidance for perturbed inputs.}
Perturb each input by a relative factor \(1+O(d_i^{-.1})\) on its
support, and begin with independent group choices. The failure probability
of a star requirement, also after any positive group-bin pin, is at most
\(2\epsilon_i^{1/8}+n^{-A_c+o(1)}\). Indeed pool typicality and Markov
control the conditional label failure, the perturbation is bounded even
across \(h_i\) groups, and
\(\max_D\widetilde q_g(D)\le n^{-A_c+o(1)}\) bounds repeated bins.
Each group belongs to only \(O(h_i^2)\) such stars.

Enforce the column bound by the capacity subset certificates of
Section~\ref{sec:bin-and-injection-stages}, using the present nominal
loads and the target \(.2\). The incoming bound \(\theta_*\) leaves
room for the fixed bucket multiplier. A group's contribution to a column
when it chooses that column's bin is at most \(d_i^{-1/2}\). The same
factorial calculation, inflating each certificate probability by 2 per
participating group, bounds the sum touching any specified group by
\(d_i\exp(-\Omega(d_i^{.4}))\), even after fixing that group's bin.
Star charges \(\epsilon_i^{1/16}\) together with these certificate
charges have total at most \(d_i^{-10}\) per group. This uses the fixed
choice of \(Q_0\): \(h_i\) is a large power of \(\log d_i\) and is
small compared with \(d_i^{.01}\). The products of neighboring avoidance
factors therefore fit the certificate inflations and the star charges.
Lemma~\ref{lem:conditional-avoidance} supplies a law avoiding all the
constraints. For \(l\) group targets, its upper comparison, including
the input perturbation, is within \(\exp(d_i^{-.05}l)\). Thus every
constructed law lies in the same \(\mathcal C\).

\paragraph{Relative singleton bounds and separation.}
Write \(\widehat q_g\) temporarily for a perturbed probability input.
Fix a group \(g\), and let \(A\) be avoidance of all constraints not
reading it. Under the product law \(D_g\) is independent of \(A\).
For every supported target \(D\), the pinned charge estimates and the
reciprocal avoidance factors give
\[
 \Pr(\text{some constraint touching }g\text{ fails}\mid A,D_g=D)
       \le d_i^{-2}.
\]
Indeed a touching event's nonneighbors remain independent of it after
this pin; avoiding its neighbors costs the reciprocal factors from
Lemma~\ref{lem:conditional-avoidance}. Summing the pinned inflated
probabilities gives the displayed bound with slack. For the upper bound,
removing the touching constraints costs at most \(\exp(O(d_i^{-10}))\),
by their total charge. Thus each marginal after full avoidance is
\((1\pm d_i^{-2})\widehat q_g(D)\).

For real prices \(c_g(D)\), put
\(\bar c_g=\sum_D c_g(D)\widetilde q_g(D)\). Multiply the original
input by
\(1+d_i^{-.1}\operatorname{sign}(c_g(D)-\bar c_g)\) and normalize.
As in Lemma~\ref{lem:near-product-injection}, its centered-price gain
is at least half of
\(d_i^{-.1}\sum_D|c_g(D)-\bar c_g|\widetilde q_g(D)\).
The two-sided singleton error is smaller than this gain. Hence no linear
functional separates the original marginal vector from the marginal
image of \(\mathcal C\). A law in \(\mathcal C\) with the exact
original marginals exists.
\end{proof}

\begin{proposition}[Calibrated role labels]\label{prop:calibrated-role-labels}
In low cluster mode, given the successful bins, there is an injective
label assignment with exact singleton marginals \(U_g^W(D)\), avoiding
every \(E_v^{\rm in}\). Its joint probability on \(l\) distinct queried
roles is at most \(\exp(d_i^{-.01}l)\) times the product of their laws,
provided at most \(h_i\) roles are queried per bin.

In bounded and low direct-bias modes, fix a typical cell pool
\(\mathcal P\) and put \(L=L_i^{\rm slot}\). Identifying singleton bins
with labels, there is an injective assignment with exact role marginals
\(\widetilde q_{g(b)}\). For any \(l\le L^{.025}\) distinct roles and
specified labels, it satisfies
\[
 \Pr(y_b\text{ is assigned at every queried }b\mid\mathcal P)
 \le \exp(L^{-.04}l)
          \prod_{b\ {\rm queried}}\widetilde q_{g(b)}(D(y_b)).
\]
In particular, this comparison holds for every \(l\le n^3\).
\end{proposition}

\begin{proof}
First work in cluster mode. Perturb each role law relatively by
\(1+O(d_i^{-.02})\). In each physical bin independently, apply
Lemma~\ref{lem:near-product-injection} to its roles and perturbed laws.
The atoms are at most \(d_i^{-.95}\) by \eqref{eq:source-13}, and the
column sums retain the lemma's slack from the \(.2\) bound. These
within-bin assignments have exact perturbed singleton marginals. Here the
independent product variables are whole injective assignments within
bins, rather than individual role labels.

A star failure has probability \(O(\epsilon_i^{1/8})\) under this
product of bin assignments. Its probability conditional on any supported
role-label pin in a touched bin is at most
\(O(d_i\epsilon_i^{1/8})\). To see the two cases, a pin outside the
star can be included in the joint upper comparison for at most
\(h_i+1\) roles and then canceled by division by its exact singleton
mass. If the pinned role is in the star, division costs \(O(d_i)\),
since every supported atom of its perturbed uniform-subset law is
\(\gtrsim1/d_i\). The comparison factors in both cases are bounded.

Each bin serves at most \(d_i\) roles and touches at most \(d_i h_i\)
star tests. Local-lemma conditioning therefore avoids all failures with
total charge at most \(d_i^{-10}\) per bin. The same pinned argument as
for group bins gives two-sided relative singleton control, now with
``touching'' meaning sharing the pinned role's bin variable. Let
\(\mathcal C'\) be the compact convex set of laws on injective,
failure-free assignments satisfying all the stated joint upper bounds
relative to the unperturbed role laws. The perturbation and avoidance
costs fit \(\exp(d_i^{-.01}l)\), so all price-directed constructions
belong to \(\mathcal C'\). Centered-price separation, with the
\(d_i^{-.02}\) gain dominating the singleton error, yields the exact
unperturbed marginals.

In direct modes the physical bins are singletons, but the label universe
for Lemma~\ref{lem:near-product-injection} is the whole cell pool of
\(L_i^{\rm slot}=n^{A_c+o(1)}\) labels. The role-specific permitted
restrictions have atoms \(O(1/L_i^{\rm slot})\), and typicality gives
column sums below \(.4\). Thus that lemma applies with
\(d=L_i^{\rm slot}\), not with \(d_i=1\). It gives the exact
\(\widetilde q_{g(b)}\) marginals and the displayed joint upper bound,
which applies through at least \(n^3\) queries since
\((L_i^{\rm slot})^{.025}=n^{5+o(1)}\). There are no internal tests.
Fix the resulting cell strategies as lookup kernels, so fresh drawings
read only the declared local data.
\end{proof}

\subsection{Fresh comparisons}
\label{sec:fresh-comparisons}

\begin{remark}[Exact marginals and fixed-scale errors]\label{rem:low-fixed-scale}
The low cluster parameters need not diverge along the counterexample
sequence. The joint comparison errors from the two calibration
propositions are controlled by the fixed large threshold; they need not
vanish as \(n\to\infty\). The singleton estimate below instead uses
exact calibration at both stages, together with the pool normalizers and
the negligible history-conditioning and permission-trim losses.
\end{remark}

\begin{lemma}[Fresh singleton comparison]\label{lem:fresh-singleton}
The cell procedures give probability priors \(\sigma_v\) valid internally
throughout each typical cell. If \(D(y)\) is the physical bin containing
\(y\), then each odd role \(b\) of patch \(i=i(b)\) has fresh law
\begin{equation}
 p_b(y\mid\mathcal P)\le (1+O(n^{-3}))
       (B_i/L_i^{\rm slot})\pi_i(y)
       \mathbf1_{\{D(y)\in\mathcal P,\ {\rm permitted}\}}.          \label{eq:source-22}
\end{equation}
\end{lemma}

\begin{proof}
At a fixed typical pool, exact calibration at both stages gives the
identity
\[
 p_b(y\mid\mathcal P)=
 \mathbb E_{W\mid\mathcal P,\,\text{slice and cell gates}}
       [\widetilde q_{g(b)}(D(y))U_{g(b)}^W(D(y))(y)].
\]
Omit the expectation in direct modes. The load-gate denominator costs
\(1+e^{-n^{c_0}}\). Pool restriction then costs
\((B_i/L_i^{\rm slot})(1+O(n^{-4}))\), retaining the presence and
permission indicators. Undoing the permission trim costs
\(1+e^{-\Omega(n)}\). Finally the internally pretrimmed marginal,
averaged over slice-conditioned data, is exactly \(\pi_i\) by the low
profile definition in Proposition~\ref{prop:balanced-patch-profiles}.
These factors prove \eqref{eq:source-22}.

For an iid uniform pool, a fixed bin belongs to the pool with probability
at most \(L_i^{\rm slot}/B_i\). Thus averaging
\eqref{eq:source-22} conditional on individual typicality gives at most
\[
 \frac{1+O(n^{-3})}{\Pr(\mathcal P\text{ is typical})}\,\pi_i(y)
       \le(1+O(n^{-3}))\pi_i(y).
\]
Both measures are probabilities, so this also gives total-variation
closeness. This is an iid-pool comparison; it does not assert independence
of the actual permutation pools.
\end{proof}

\begin{lemma}[Fresh internal-prior comparison]\label{lem:fresh-internal-prior}
On unforced iid slot pools, with individual typicality retained, the law
of a fresh internal prior has bounded upper comparison with its
unrestricted base experiment. More precisely, for every nonnegative test
\(\Phi\) of the prior, with \(\Phi(0)=0\),
\[
 \mathbb E_{\mathcal P\ {\rm iid}}
  \left[\mathbf1_{\{\mathcal P\ {\rm typical}\}}
          \mathbb E_{\rm fresh(\mathcal P)}\Phi(\sigma_v)\right]
       \le K\,\mathbb E_{\rm base}\Phi(\sigma_v).
\]
Consequently its mean and its permission-exception estimates hold up to
fixed factors. True internal validity may be retained in the comparison.
\end{lemma}

\begin{proof}
Remove the cell-load conditioning at its negligible upper cost, keeping
its gate until the conditional samplers have been compared. Apply first
the label upper comparison and then the group-bin upper comparison on the
internal star. The queried groups choose distinct physical bins; retain
that restriction while averaging the pool. For each target bin replace
its \(\widetilde q_g\) factor by the upper bound consisting of
\(B_i/L_i^{\rm slot}\), its pool-presence indicator, and the raw
\(q_g^W\) factor, with the trim and normalizer costs.

If there are \(m\) distinct target bins, simultaneous containment in
unforced iid slots has probability at most
\((L_i^{\rm slot}/B_i)^m\): sum over distinct slot indices realizing
those targets. This cancels all \(m\) restriction factors. The remaining
nonnegative integral can now discard distinctness and use the raw bin and
reference label kernels. Removing slice conditioning on
\eqref{eq:source-15} restores the unrestricted base experiment. The raw
posterior rule itself has stayed unchanged throughout.

All costs together are bounded: besides the normalizers, the relevant
quantities are \(h_i^{O(1)}\sqrt{\epsilon_i}\ll1\) for the pretrims and
\(h_i d_i^{-.01}\ll1\) for the joint comparisons. In direct modes the
internal prior is deterministic. The same proof can retain the true
internal-validity indicator. Taking \(\Phi(\sigma)=\sigma(x)\) gives
the mean estimate from \eqref{eq:source-14}; taking the permission
exception indicator gives \eqref{eq:source-21} up to the same fixed
factor. Conditioning the iid pool on typicality costs its probability's
reciprocal. An own-pool slot pin is excluded here because it would change
the containment calculation.
\end{proof}
